\xchapter{Dedication}




TO
THE


MOST
REVEREND FATHER IN GOD


WILLIAM


LORD ARCHBISHOP
OF CANTERBURY,


PRIMATE OF ALL
ENGLAND,


FORMERLY
REGIUS PROFESSOR OF DIVINITY IN THE UNIVERSITY OF OXFORD,


THIS LIBRARY


OF


ANCIENT
BISHOPS, FATHERS, DOCTORS, MARTYRS, CONFESSORS,


OF CHRIST’S
HOLY CATHOLIC CHURCH,


IS


WITH HIS
GRACE’S PERMISSION


RESPECTFULLY
INSCRIBED,


IN
TOKEN OF


REVERENCE FOR
HIS PERSON AND SACRED OFFICE,


AND
OF


GRATITUDE FOR
HIS EPISCOPAL KINDNESS.

\xchapter{The Catechetical Lectures of S. Cyril, Archbishop of Jerusalem, Translated}

[4th Edition, Parker \& Co. and
Rivingtons, 1872]

\xsection{Preface}

S. CYRIL,
the author of the Catechetical Lectures which follow, was born in an
age ill adapted for the comfort or satisfaction of persons
distinguished by his peculiar character of mind, and in consequence
did not receive that justice from contemporaries which the Church
Catholic has since rendered to his memory. The Churches of\emph{}Palestine,
apparently his native country, were the first to give reception to
Arius on his expulsion from Alexandria, and without adopting his
heresy, affected to mediate and hold the balance between him and his
accusers. They were followed in this line of conduct by the provinces
of Syria and Asia Minor, till the whole of the East, as far as it was
Grecian, became more or less a large party, enduring to be headed by
men who went the whole length of Arianism, from a fear of being
considered Alexandrians or Athanasians, and a notion, for one reason
or other, that it was thus pursuing a moderate course, and avoiding
extremes. What were the motives which led to this perverted view of
its duty to Catholic truths then so seriously endangered, and what the
palliations in the case of individuals, need not be minutely
considered here. Suffice it to say, that between the Churches of Asia
and the metropolis of Egypt there had been distinctions, not to say
differences and jealousies of long standing; to which was added this
great and real difficulty, that a Council held at Antioch about sixty
years before had condemned the very term, Homoüsion, which was the
symbol received at Nicæa, and maintained by the Alexandrians. The
latter were in close agreement with the Latin Church, especially
with Rome; and thus two great confederacies, as they may be called,
were matured at this distressing era, which outlived the controversy
forming them, the Roman, including the West and Egypt, and the
Asiatic, extending from Constantinople to Jerusalem. Of the Roman
party, viewed at and after the Arian period, were Alexander,
Athanasius, Eustathius, Marcellus, Julius, Ambrose, and Jerome; of the
Asiatic, Eusebius of Cæsarea, Cyril, Meletius, Eusebius of Samosata,
Basil of Cæsarea (the Great), Basil of Ancyra, Eustathius of Sebaste,
and Flavian. Of the latter, some were Semi-arian; of the former, one
at least was Sabellian; while the majority of both were, to say the
least, strictly orthodox; some of the latter indeed acquiescing with
more or less of cordiality in the expediency of adopting the important
Symbol of the Nicene Council, but others, it need scarcely be said, on
both sides, being pillars of the Church in their day, as they have
been her lights since. Such was the general position of the Church;
and it is only confessing that the early Bishops and Divines were men ``of
like passions with'' ourselves, to add, that some of them sometimes
misunderstood or were prejudiced against others, and have left on
record reports, for the truth of which they trusted perhaps too much
to their antecedent persuasions, or the representations of their own
friends. When Arianism ceased to be supported by the civil power, the
controversy between East and West died; and peace was easily effected.
And the terms of effecting it were these:—the reception of the Homoüsion
by the Asiatics, and on that reception their recognition, in spite of
their past scruples, by the Alexandrians and Latins. In this sketch
the main outlines of S. Cyril's history will be found to be contained;
he seems to have been afraid of the term Homoüsion \footnote{1. v. Bened. note iv. 7. xvi.
23.}, to have been disinclined both to the friends of Athanasius and
to the Arians \footnote{2. Lect.
iv. 8. xi. 12, 16, 17. xv. 9.}, to have
allowed the tyranny of the latter, to have shared in the general
reconciliation, and at length both in life and death to have received
honours from the Church, which, in spite of whatever objections
may be made to them, appear, on a closer examination of his history,
not to be undeserved.

CYRIL is said to have
been the son of Christian parents, but the date and place of his birth
is unknown. He was born in the first years of the fourth century, and
at least was brought up in Jerusalem. He was ordained Deacon probably
by Macarius, and Priest by Maximus, the Bishops of Jerusalem; the
latter of whom he succeeded A.D.
349, or 350. Shortly before this, (A.D.
347, or 348,) during his Priesthood, he had delivered the Catechetical
Lectures which have come down to us. With his Episcopate commence the
historical difficulties under which his memory labours. It can
scarcely be doubted that one of his consecrators was Acacius of Cæsarea
\footnote{3. v. Diss.
Beded. P. xviii. sq.}, the leader of Arianism
in the East, who had just before (A.D. 347) been
deposed by the Council of Sardica; yet, as the after history shows,
Cyril was no friend of the Arians or of Acacius \footnote{4. Theod.
ii. 26.}. He was canonically consecrated by the Bishops of his province,
and as Acacius was still in possession of the principal see, he was
compelled to a recognition which he might have wished to dispense
with. He seems to have been a lover of peace; the Council of Sardica
was at first as little acknowledged by his own party as by the Arians;
and Acacius, being even beyond other Arians skilful and subtle in
argument, and admitting the special formula \footnote{5. The [\emph{kata
panta homoion}]. Vid. Lecture iv. 7. xi. 4, 9, 18.} of Cyril on the doctrine in controversy, probably succeeded in
disguising his heresy from him.

A more painful account,
however, of his consecration is given by S. Jerome \footnote{6. Jer.
Chron. Socr. ii. 38. Sozom. iv. 20.}, supported in the main by other writers, which can only be
explained by supposing that Father to be misled by the information or
involved in the prejudices of Cyril's enemies. He relates, that upon
Maximus's death, the Arians seized upon the Church of Jerusalem,
and promised Cyril the see on condition of his renouncing the
ordination he had received from Maximus, and submitting to
re-ordination from their hands; that he assented, served in the Church
as a mere Deacon, and was then raised by Acacius to the Episcopate,
when he persecuted Heraclius, whom Maximus had consecrated as his
successor. This account, incredible in itself, is contradicted, on the
one hand, by the second General Council, which in its Synodal letter
plainly states that he had been ``canonically ordained'' Bishop, and on
the other by his own writings, which as plainly show, that in doctrine
he was in no respect an Arian or an Arianizer.

If he suffers in memory from
the Latin party as if Arian, he suffered not less in his life from the
Arians as being orthodox. Seven or eight years after his consecration,
he had a dispute with Acacius about the rights of their respective
Churches \footnote{7. Socrat.
ii. 40. Sozom. iv. 25. Theod. ii. 27.}. Acacius in
consequence accused him to the Emperor Constantius of holding with the
orthodox; to which it was added that he had during a scarcity sold
some offerings made by Constantine to his Church, to supply the wants
of the poor. Cyril in consequence was deposed, and retired to Tarsus \footnote{8. Theodor.
ii. 26.}; where, in spite of the efforts of Acacius, he was hospitably
received, and employed by Silvanus the Semi-arian Bishop of the place.
We find him at the same time in friendship with Eustathius of Sebaste
and Basil of Ancyra, both Semi-arians \footnote{9. Sozom.
iv. 25. Philostorg. iv. 12.}. His own writings, however, as has already been intimated, are
most exactly orthodox, though he does not in the Catechetical
Lectures, use the word Homoüsion; and in associating with these men
he went little farther than S. Hilary \footnote{10. Hilar.
de Synod. 77, 88. \&c. v. fragm. II. 4. (Ed. Ben. Cyr. p. lx. D.)}, during his banishment in Asia Minor, who calls Basil and
Eustathius ``most holy men,'' than S. Athanasius, who acknowledges as ``brethren''
those who but scrupled at the \emph{word} Homoüsion \footnote{11. Athan.
de Synod. 41.}, or than S. Basil of Cæsarea, who till a late period of
his life was an intimate friend of Eustathius.

In A.D.
359, two years after his deposition, he successfully appealed against
Acacius to the Council of Seleucia \footnote{12. v.
Dissert. 1. Ed. Bened. p. lvi. F.}, one of the two branches of the great Council of East and
West, which was convened under the patronage of Constantius to settle
the troubles of the Christian world. But the next year, Acacius
contriving to bring the matter before a Council at Constantinople,
where the Emperor was staying, Cyril with his friends was a second
time deposed, and banished from Palestine \footnote{13. Sozom.
iv. 25. Philostorg. iv. 12. v. 1.}.

On Constantius's death all the
banished bishops were restored \footnote{14. Sozom.
v. 5.}, and Cyril, who was at that time with Meletius of Antioch,
returned to Jerusalem, A.D. 362. He was
there at the time of Julian's attempt to rebuild the Temple \footnote{15.
Socrat. iii. 20. Ruffin. i. 37.}, and from the Prophecies boldly foretold its failure.

He was once more driven from
his see, during the reign of the Arian Valens \footnote{16. Sozom.
iv. 30. Jerom. Script. Eccles. 112.}, (A.D.
367,) and he remained dispossessed till A.D.
378.

About the time of the death of
Valens, the last of the Arian princes, he was restored, but under what
circumstances is unknown. The Arians fell once for all with their
imperial protectors; and soon after, that union of Christian Churches
took place, which would never have been interrupted, had not a few
bold and subtle-minded men contrived to delude them into the belief of
mutual differences. S. Athanasius, the great peace-maker of the
Church, was gone to his rest; and S. Basil also, who had mourned over
evils which he had no means of remedying. Gregory of Nyssa, the
brother of the latter, Gregory of Nazianzum, Meletius of Antioch,
remained; and were present together with Cyril in the second General
Council \footnote{17.
Socrat. v. 8. Sozom. vii. 7. v. Dissert. Bened. p. lxxxii.}, which
formally restored the latter to his see, and in its letter to the
Western Bishops speaks of him as ``the most reverend and religious
Cyril, long since canonically appointed by the Bishops of the
province, and in many ways and places a withstander of the Arians \footnote{18. Theod.
Hist. v. 9.}.'' He died about the year 386. Except one or two short
compositions and fragments, nothing remains of his writings but his
Catechetical Lectures.

No ecclesiastical writer could
be selected more suitable to illustrate the main principle on which
the present ``Library of the Fathers'' has been undertaken, than S.
Cyril of Jerusalem. His Catechetical Lectures were delivered, as we
have seen, when he was a young man; and he belonged, till many years
after their delivery, to a party or school of theology, distinct, to
say the least, from that to which the most illustrious divines of his
day belonged; a school, never dominant in the Church, and expiring
with his age. It is not then on the score of especial personal
authority that his Lectures are now presented to the English reader;
and if the simple object of this Publication were to introduce the
latter to the wise\emph{}and good of former times, S. Cyril would
have no claims to a place in it beyond many who have lived since.

In saying this, it is far
indeed from being asserted, that the personal claims of the Fathers of
the Church on our deference are inconsiderable; for it happens, not
unnaturally, that the works which have been preserved, were worth
preserving, or rather that their writers would have been extraordinary
men in any age, and speak with the weight of great experience,
ability, and sanctity. To those who believe that moral truth is not
gained by the mere exercise of the intellect, but is granted to moral
attainments, and that God speaks to inquirers after truth by the mouth
of those who possess these, the writings of S. Basil or S. Augustine
must always have an authority independent of their date or their
agreement; nor is it possible for serious persons to read them,
without feeling the authority which they possess as individuals. This,
however, whatever it be, is not the main subject to which the
present Translations propose to direct attention. The works to be
translated have been viewed simply and plainly in the light of
witnesses to an historical fact, viz. the religion which the Apostles
transmitted to the early Churches, a fact to be ascertained as other
past facts, by testimony, requiring the same kind of evidence, moral
not demonstrative, open to the same difficulties of proof, and to be
determined by the same practical judgment. It seems hardly conceivable
that a fact so public and so great as the religion of the first
Christians should be incapable of ascertainment, at least in its
outlines, that it should have so passed away like a dream, that the
most opposite opinions may at this day be maintained about it without
possibility of contradiction. If it was soon corrupted or
extinguished, then it is obvious to inquire after the history of such
corruption or extinction; such a revolution every where, without
historical record, being as unaccountable as the disappearance of the
original religion for which it is brought to account. At first sight
there is, to say the least, a considerable antecedent improbability in
the notion, that, whereas we know the tenets and the history of the
Stoic or the Academic philosophy, yet we do not know the main tenets,
nor yet the fundamental principles, nor even the spirit and temper of
Apostolic Christianity.

Under a sense of this
improbability, in other words with an expectation that historical
research would supply what they sought, our Divines at and since the
Reformation have betaken themselves to the extant documents of the
early Church, in order to determine thereby what the system of
Primitive Christianity was; and so to elicit from Scripture more
completely and accurately that revealed truth, which, though revealed
there, is not on its surface, but needs to be \emph{deduced}and \emph{developed}from it. They went to the Fathers for information concerning
matters, on which the Fathers at first sight certainly do promise to
give information, just as inquirers into any other branch of knowledge
might study those authors who have treated of it; and, whether
or not they found what they sought, it surely was reasonable so to
seek it, and cannot be condemned except by the \emph{event},—that\emph{}is, by showing that their expectation, however reasonable
antecedently, is mistaken in fact, that after the search into history,
no evidence is forthcoming concerning the tenets, nor yet the
principles, nor even the temper which the Apostles inculcated.

A like expectation has
actuated the present Publication; it has been conceived probable, to
say the least, that the study of the writings of the Fathers will
enable us to determine morally, to make up our minds for practical
purposes, what the \emph{doctrines}of the Apostles were, for instance
whether or not they believed in our Lord's Divinity, or the general
necessity of Baptism for salvation;—or if not the doctrines, still
what were their \emph{principles}, as whether or not or how far they
allowed of using secular means for advancing Christian truth, or
whether or not they sanctioned the monarchical principle, or again the
centralizing principle, or again the principle of perpetuity in Church
matters, or whether they considered that Scripture should be
interpreted in the mere letter, or what is called spiritually;—or at
least what was their \emph{temper}, for instance, whether or not it
was what is now called in reproach superstitious, or whether or not
exclusive, or whether or not opposed to display and excitement. On
some or other of these points there are surely grounds for expecting \emph{information}
from the Fathers, sufficient for our practice, and therefore having
claims upon it.

Recourse then being had to the
writings of the Fathers, in order to obtain information as to this
historical fact, viz. the doctrines, the principles, and the religious
temper of Apostolical Christianity, so far we have little to do with
the personal endowments of the Fathers, except as these bear on the
question of their fidelity. Their being men of strictest lives and
most surpassing holiness, would not prove that they knew what it was
the Apostles taught; and, were they but ordinary men, this need
not incapacitate them from being faithful witnesses and serviceable
informants, if they were in a position to be such. We should have only
to take into account, and weigh against each other, their
qualifications and disqualifications, for being evidence to a fact; we
should have to balance honesty against prejudice, education against
party influence, early attachments against reason, and so on. Thus we
should treat them, taken one by one; but even this sort of personal
scrutiny will be practically superseded, when we consult them, not
separately, but as our Reformed Church ever has done, together; and
demand their \emph{unanimous} testimony to any point of doctrine or
discipline, before we make any serious use of them: for it stands to
reason, that, where they agree, the peculiarities of their respective
nations, education, history, and period, instead of suggesting an
indefinite suspicion against the subject-matter of their testimony,
does but increase the evidence of its truth. Their testimony becomes
the concurrence of many independent witnesses in behalf of the same
facts; and, if it is to be slighted or disparaged, one does not see
what knowledge of the past remains to us, or what matter for the
historian. Viewing S. Cyril, for instance, as one of a body who bears
a concordant evidence to the historical fact, that the Apostles taught
that Christ is God \footnote{19. iv.
7, \&c.}, or
that Baptism is the remedy of original sin \footnote{20.
Introd. 16. ii. 5. iii.4, 11. xii. 15.}, or that celibacy is not imperative on the clergy \footnote{21. xii.
25.}, whether he was Asiatic or African, of the Roman or the
Oriental party, as little matters, as when we consider him as one of a
company bearing witness to the historical fact, that the Apostles and
their associates wrote the New Testament. Indeed, as the matter
stands, there is something very remarkable and even startling to the
reader of S. Cyril, to find in a divine of his school such a perfect
agreement, for instance as regards the doctrine of the Trinity, with
those Fathers who in his age were more famous as champions of it. Here
is a writer, separated by whatever cause from what, speaking
historically, may be called the Athanasian School, suspicious of its
adherents, and suspected by them; yet he, when he comes to explain
himself \footnote{22. iv.
7. vi. 1. x. and xi. xii. 1.}, expresses
precisely the same doctrine as that of Athanasius or Gregory, while he
merely abstains from the particular theological term in which the
latter Fathers agreeably to the Nicene Council conveyed it. Can we
have a clearer proof that the difference of opinion between them was
not one of ecclesiastical and traditionary doctrine, but of practical
judgment? that the Fathers at Nicæa wisely considered that, under the
circumstances, the word in question was the only symbol which would
secure the Church against the insidious heresy which was assailing it,
while S. Cyril \footnote{23. v.
Hilar. Contr. Const. 3. 12. Athan. de Synod. 12. (ed. Bened. Cyr. lii.
B.)}, with
Eusebius of Cæsarea, Meletius, and others, shrunk from it, at least
for a while, as if an addition to the Creed, or a word already taken
into the service of an opposite heresy \footnote{24. vid.
Bull. Defens. F. N. ii. 1.}, and likely to introduce into the Church heretical notions?
Their judgment, which was erroneous, was their own; their faith was
not theirs only, but shared with them by the whole Christian world.

At the same time it must be
granted, that this view of the Fathers as witnesses to Apostolic truth
not individually but collectively, clear and unanswerable as it is,
considered as a view, is open to some great practical inconveniences,
when acted on in such an undertaking as that in which the present
Editors are engaged. For since, by the supposition, no one of the
Fathers is necessarily right in all his doctrine, taken by himself,
but may be erroneous in secondary points, each taken by himself is in
danger, by his own peculiarities, on the one hand of throwing
discredit on all together, on the other of perplexing those who by
means of the Fathers are inquiring after Catholic truth. And whereas
in any publication of this nature, they cannot appear all at once, but
first one and then another, and at all events cannot be read
altogether, it follows that, during their gradual perusal, unavoidable
prejudice will often attach to the Fathers, and to the Catholic Faith,
and to those who are enforcing the latter by means of the former. And
thus Editors of the Fathers are pretty much in the condition of
Architects, who lie under the disadvantage, from which Painters and
Sculptors are exempt, of having their work exposed to public criticism
through every stage of its execution, and being expected to provide
symmetry and congruity in its parts independent of the whole.

Such are the circumstances in
which we find ourselves, open to remark for every opinion, every
sentence, every phrase, of every Father, before its meaning,
relevancy, importance, or bearings are ascertained; before it is known
whether it will be, as it were, obliterated by others, or completed,
or explained, or modified, or unanimously witnessed. And since the
evil is in the nature of the case itself, we can do no more than have
patience, and recommend patience to others, and with the racer in the
Tragedy look forward steadily and hopefully to the \emph{event},—[\emph{T\={O}i
TELEI pistin pher\={o}n}],—when,\emph{}as we trust, all that
is inharmonious and anomalous in the details, will at length be
practically smoothed. Meanwhile, as regards the condition of the
reader himself, we consider that we shall sufficiently provide for his
perplexity by reminding him of his duty to take his own Church for the
present as his guide, and her decisions as a key and final arbiter, as
regards the particular statements of the separate Fathers, which he
may meet with; being fully confident, that her judgment which he
begins by taking as a touchstone of each, will in the event be found
to be really formed, as it ought to be, on a view of the testimony of
all.

In expressing, however, these
thoughts, it is obvious to anticipate an objection of another sort
which is likely to be urged against our undertaking, to the effect
that all these dangers and warnings are gratuitous, Scripture itself
containing sufficient information concerning the doctrines,
principles, and mind of the Apostles, without having recourse to the
difficult and, as has been above confessed, the anxious task, whether
ultimately successful or not, of collecting the object of our inquiry
from the writings of the Fathers. This is not the place to treat of an
objection, to which much attention has been drawn for several years
past; yet thus much may be observed in passing:—If the sufficiency
of Scripture for teaching as well as proving the Christian faith be
maintained as a theological truth, the grounds in reason must be
demanded, grounds such as are independent of that inquiry into history
which it is brought forward to prohibit. If it is urged as a truth
obvious in matter of fact, and practically certain, then its
maintainers have to account for the actual disagreement among readers
of Scripture as to \emph{what} the faith, principles, and temper of
the Apostles were. And if it be urged on the authority of the sixth
Article of our Church, then they must be asked, why if this Article
contained a reason against deferring to Antiquity, the Convocation of
1571, which imposed it, at the same time, as is well known, ordered
all preachers to teach \emph{according to the}Catholic Fathers, and
why our most eminent Divines, beginning with the writers of the
Homilies themselves, have ever pursued that very method.

Nothing can be more certain
than that Scripture contains all necessary doctrine; yet nothing, it
is presumed, can be more certain either, than that, practically
speaking, it needs an interpreter; nothing more certain than that our
Church and her Divines assign the witness of the early ages of
Christianity concerning Apostolic doctrine, as that interpreter.

Without, however, entering
into a question which our Church seems to have determined for us, a
few words shall be devoted to the explanation of a verbal difficulty
by which it is often perplexed. An objection is made, which, when
analyzed, resolves itself into the following form. ``Either
Antiquity does or does not teach something over and above Scripture:
if it does, it adds to the inspired word;—if it does not, it is
useless.—Does it then or does it not \emph{add} to Scripture?'' And,
as if showing that the question is a perplexed one, of various writers
who advocate the use of Antiquity, one may be found to speak of the
writings of the Fathers as enabling us to ascertain and revive truths
which have fallen into desuetude, while another may strenuously
maintain, that they impart the knowledge of no new truths over and
above what Scripture sets before us. Now, not to touch upon other
points suggested by this question, it may be asked by way of
explanation, whether the exposition of the true sense of any legal
document, any statute or deed, which has been contested, is an
addition to it or not? It is in point of \emph{words} certainly; for
if the words were the same, it would be no explanation; but it is no
addition to the sense, for it professes to be neither more nor less
than the very sense, which is expressed in one set of words in the
original document, in another in the comment. In like manner, when our
Saviour says, ``I and My Father are One,'' and Antiquity interprets ``One''
to mean one in substance,'' this \emph{is} an addition to the \emph{wording},
but no addition to the sense. Of many possible means of interpreting a
word, it cuts away all but one, or if it recognizes others, it reduces
them to harmony and subordination to that one. Unless the Evangelist
wished his readers to be allowed to put any conceivable sense upon the
word, the power of doing so is no privilege; rather it is a privilege
to know that very meaning, which to the exclusion of all others is the
true meaning. Catholic Tradition professes to do for Scripture just
which is desirable, whether it is possible or not, to relieve us from
the chance of taking one or other of the many senses which are wrong
or insufficient, instead of the one sense which is true and complete.

But again, every diligent
reader of the Bible has a certain idea in his own mind of \emph{what}
its teaching is, an idea which he cannot say is gained from this or
that particular passage, but which he has gained from it as a whole,
and which if he attempted to prove argumentatively, he might perplex
himself or fall into inconsistencies, because he has never trained his
mind in such logical processes; yet nevertheless he has in matter of
fact a \emph{view} of Scripture doctrine, and that gained from
Scripture, and which, if he states it, he does not necessarily state
in words of Scripture, and which, whether after all correct or not, is
not incorrect merely because he does not express it in Scripture
words, or because he cannot tell whence he got it, or logically refer
it to, or prove it from, particular passages. One man is a Calvinist,
another an Arminian, another a Latitudinarian; not logically merely,
but from the impression gained from Scripture. Is the Latitudinarian
necessarily adding to Scripture because he maintains the \emph{proposition},
``religious opinions matter not, so that a man is sincere,'' a
proposition not \emph{in terminis} in Scripture? Surely he is
unscriptural, not because he uses words not in Scripture, but because
he thereby expresses ideas which are not expressed in Scripture. In
answer then to the question, whether the Catholic system is an
addition to Scripture, we reply, in one sense it is, in another it is
not. It is not, inasmuch as it is not an addition to the range of
independent \emph{ideas} which Almighty God intended should be
expressed and conveyed on the whole by the inspired text: it is an
addition, inasmuch as it is in addition to their \emph{arrangement},
and to the \emph{words} containing them,—inasmuch as it stands as a
conclusion contrasted with its premisses, inasmuch as it does that
which every reader of the Scriptures does for himself, express and
convey the ideas more explicitly and determinately than he finds them,
and inasmuch as there may be difficulty in duly referring every part
of the explicit doctrine to the various parts of Scripture which
contain it.

Nothing here is intended
beyond setting right an ambiguity of speech which both perplexes
persons, and leads them to think that they differ from others, from
whom they do not differ. No member of the English Church ever thought
that the Church's creed was an \emph{addition} to Scripture in any
other sense than that in which an individual's own impression
concerning the sense of Scripture is an addition to it; or ever
referred to a supposed deposit of faith distinct from Scripture
existing in the writings of the Fathers, in any other sense than in
that in which asking a friend's opinion about the sense of Scripture,
might be called imputing to him unscriptural opinions. The question of
words then may easily be cleared up, though it often becomes a
difficulty; the \emph{real} subject in dispute, which is not here to
be discussed, being this, \emph{how} this one true sense of Scripture
is to be learned, whether by philological criticism upon definite
texts,—or by a promised superintendence of the Holy Ghost teaching
the mind the true doctrines from Scripture, (whether by a general
impression upon the mind, or by leading it, text by text piecemeal
into doctrine by doctrine;)—or, on the other hand, by a blessing of
the Spirit upon studying it in the right way, that is, in the way
actually provided, in other words, according to the Church's
interpretations. In all cases the text of Scripture and an exposition
of it are supposed; in the one the exposition comes first and is
brought to Scripture, in the other it is brought out after examination
into Scripture; but you cannot help assigning some exposition or
other, if you value the Bible at all. Those alone will be content to
ascribe no sense to Scripture, who think it matters not whether it has
any sense or not. As to the case of a difference eventually occurring
in any instance of importance, between what an individual considers to
be the sense of Scripture and that which he finds Antiquity to put
upon it, the previous question must be asked, whether such difference
is likely to arise. It will not arise in the case of the majority, nor
in the case of serious, sensible, and humble minds; and where
men are not such, it will be but one out of many difficulties. A
person however, thus circumstanced, whether from his own fault or not,
\emph{is} in a difficulty; difficulties are often our lot, and we must
bear them, as we think God would have us. We can cut the knot by
throwing off the authority of the Fathers; and we can remain under the
burden of the difficulty by allowing that authority; but, however we
act, we have no licence to please our taste or humour, but we act
under a responsibility.

Two main respects have been
mentioned, in which the concordant testimony of the Fathers may be
considered to throw light upon the sense of Scripture: on these a few
words are now necessary with a special reference to S. Cyril,—first
as regards the \emph{doctrine} of Scripture, next as regards the \emph{interpretation
of texts}. Now it will be found that they are more concordant as to
the doctrines themselves contained in Scripture, than as to the
passages in which these are contained and their respective force; and,
again, that they are more concordant in their view of the principles
upon which Scripture is to be interpreted, than in their application
of these principles, and their view of the sense in consequence to be
assigned to particular texts. This was to be expected, as may easily
be made appear.

There seems to have been no
Catholic exposition of Scripture, no traditionary comment upon its
continuous text. The subject-matter of Catholic tradition, as
preserved in the writings of the Fathers, is, not Scripture \emph{interpretation}
or \emph{proof}, but certain \emph{doctrines}, professing to be those
of the Gospel: and since among these we find this, ``that Scripture
contains all the Gospel doctrines,'' we infer, that, according to the
mind of the Fathers, those very doctrines which they declare to be the
Christian faith are contained in and are to be proved from Scripture.
But \emph{where} they occur in Scripture cannot be ascertained from
the Fathers, except so far as the accidental course of controversy has
brought out their joint witness concerning certain great
passages, on which they do seem to have had traditionary information.
The Arian and other heresies obliged them to appeal to Scripture in
behalf of a certain cardinal doctrine which they held by uninterrupted
tradition; and thus have been the means of pointing out to us
particular texts in which are contained the great truths which were
assailed. But while we are thus furnished with a portion of the
Scripture proof of Catholic doctrine, guaranteed to us by the
unanimous consent of the Church, it is natural also, under the
circumstances above mentioned, that many of the discussions which
occurred should contain appeals to Scripture of a less cogent
character, and evidencing the exercise of mere private judgment upon
the text in default of Catholic Tradition. The early Church had read
Scripture not for argument but for edification; it is not wonderful
that though holding the truth, and seeing it in the inspired text, and
often seeing there what we fail to see, she should nevertheless be as
little able to distribute exactly each portion of the truth to each of
its places in the text, and to analyze the grounds of those
impressions which the whole conveyed, as religious persons in the
private walks of life may be now-a-days. Accordingly her divines, one
by one, while they witness to the truth itself most sufficiently, as
speaking from Tradition, yet often prove it insufficiently, as relying
necessarily on private judgment.

For instance, the text, \emph{He
that hath seen Me, hath seem the Father}, is taken by S. Cyril,
agreeably with other early writers, as a proof that Christ is in all
things like ([\emph{homoios en pasin}]) to the Father; (Lect. xi. 18.)
and the text, \emph{Except a man be born of water and of the Spirit},
as a proof of the necessity of Baptism. (Lect. iii. 4.)\emph{}But
though there are many of equal cogency, there are many also, about
which there may be fairly difference of opinion, as when he
interprets, \emph{Surely God is in thee} (Isa. 45, 14) of the
Indwelling of the Father in the Son. (Lect. xi. 16.)

And, while it is not at all
surprising even though the Fathers should occasionally adduce texts as
proofs of certain doctrines which are not so, neither is it strange
that they should overlook proofs which did exist, and which we are
able to discern. For they were in the light of a recent Tradition; we
are in the twilight of a distant age; and our minds, like eyes
accustomed to the twilight, may discern much in the dark parts of
Scripture, which were hid from them by their very privilege.

Such imperfections, however,
in the Scripture proofs adduced by the Fathers, whether in excess or
defect, do not interfere at all with their maintenance of the great
principles that there is a Faith, and that it is in Scripture. As far
as S. Cyril is concerned, the following passages witness both truths
clearly. ``This Seal,'' he says, speaking of the Creed, ``have thou ever
in mind; which now by way of summary has been touched on in its heads,
and, if the Lord grant, shall hereafter be set forth according to our
power with Scripture proofs. For concerning the divine and sacred
mysteries of the Faith, we ought not to deliver even the most casual
remark without the Holy Scriptures; nor be drawn aside by mere
probabilities, and the artifices of argument. Do not then believe me
because I tell you these things, unless thou receive from the Holy
Scriptures the proof of what is set forth; for this salvation, which
is of our faith, is not by ingenious reasonings, but by proof from the
Holy Scriptures.'' (Lect. iv. 17.)

Again: ``Take thou and hold
that faith only as a learner and in profession, which is by the Church
delivered to thee, and is established from all Scripture. For since
all cannot read the Scripture, but some as being unlearned, others by
business, are hindered from the knowledge of them, in order that the
soul may not perish for lack of instruction, in the Articles which are
few we comprehend the whole doctrine of the faith ... Commit to memory
the Faith, merely listening to the words; and expect at the fitting
season the proof of each of its parts from the Divine
Scriptures. For the Articles of the Faith were not composed at the
good pleasure of man; but the most important points chosen from all
Scripture, make up the one teaching of the Faith. And as the mustard
seed in a little grain contains many branches, thus also this Faith,
in a few words, has enfolded in its bosom the whole knowledge of
godliness contained both in the Old and New Testaments.'' (Lect. v.
12.) The doctrine, expressed in these and other passages of S. Cyril,
is implied and assumed in a most striking way in a number of others \footnote{25. xi.
12. xii. 5. xiii. 8, 9. xiv. 2. xvi. 1, 2, 24.}.

So much on the Scripture proof
of doctrine as contained in the Fathers; as to the doctrinal sense of
Scripture, the second point to be spoken of, what has been already
observed is quite consistent, not to say connected with the remark to
be made concerning it, viz. that the Fathers are far more concordant
in assigning principles of Scripture interpretation, than in the
interpretation of particular passages. Indeed the very view they took
of the Bible led to variety, apparent discordance, and private
conjecture in interpreting it. They considered it to be a sort of
storehouse of sacred treasures \footnote{26. xi.
12.}, contained under the letter in endless profusion, piled, as it
were, one on another, with order indeed and by rule, but still often
so deeply lodged within the text, that from ordinary eyes they were
almost hidden \footnote{27. xii.
16. xiii. 14. ix. 13.}. Hence
it was considered as a duty and privilege proposed to the Christian,
to find out the ``wondrous things of God's law,'' and no meaning was so
remote from the literal text as to be proved thereby to be foreign to
it in the Divine intention. While then, according to their disposition
or school of theology, they were led, more or less, to attempt to
search into the deep mysteries of Scripture for themselves \footnote{28. iii.
16. vi. 28, 29. xii. 19.}, they felt little difficulty in multiplicity of
interpretations, little fear of inconsistency. And while such a
principle as has been described necessarily led them to diversity
in their interpretations, that diversity does but increase our
evidence of the fact of their one and all holding that principle; and
thus, while their value as commentators varies with their personal
qualifications, their adherence to that principle comes to us as a
Catholic tradition.

Instances of individual,
local, or transitory opinion, that is, of what would at present,
rightly or wrongly, be called fancifulness and caprice, are frequent
in S. Cyril's Lectures, and scarcely need specifying. Such, for
example, is his interpreting, ``Look unto the rock whence ye are hewn,''
of the Holy Sepulchre, (Lect. xiii. 35.) or ``At evening time it shall
be light,'' of the circumstances of the Crucifixion; (Lect. xiii. 24.)
and much more his considering, ``Thou hast wrought salvation in the \emph{midst
of the earth},'' (Lect. xiii. 28.) to allude to Golgotha, and ``the
fountain sealed,'' to Christ in the sepulchre after the sealing of the
stone. (Lect. xiv. 5.)

These interpretations, whether
his own or not, and whether true or not, do not profess to be
traditional, and are but witnesses, to the great principle from which
they proceed, of the everliving intelligence, deep and varied meaning,
and inexhaustible fulness of Holy Scripture. This indeed he himself
declares in one place in words which may be suitably extracted. After
giving two conjectures concerning the doctrinal meaning of the Blood
and Water, which came from our Lord's side, viz. that it typifies the
Jews' imprecation of His blood upon them and Pilate's washing his
hands of it, or again the condemnation of the Jews and the baptismal
pardon of Christians, he adds, ``\emph{For nothing happened without a
meaning}, ([\emph{ouden eik\={e} gegonen}].) Our fathers who
have written comments, have given another reason of this matter. For
since in the Gospel the power of salutary Baptism is two-fold, that
bestowed by means of water on the Illuminated, and that to holy
Martyrs in persecutions through their own blood, there came out of
that salutary side blood and water,'' \&c. (Lect. xiii. 21.)

When then it is inquired, what
information is given us by the Fathers, concerning Scripture or
Catholic doctrine, we reply, that they rather declare doctrine and say
that it is in Scripture, than prove it by Scripture, at once
concordantly and in detail; and again, that they rather tell us how we
must set about interpreting Scripture, than authoritatively interpret
it for us. It is presumed that this is on the whole correct; true as
it also is, that on a number of the most important points of doctrine
they have preserved to us, with an unanimity which is an evidence of
its Apostolic origin, the very texts in which they are contained.
Still after all the Fathers are rather led to dwell on Scripture by
itself, and on the doctrinal system by itself, as two distinct,
parallel, and substantive sources of divine information, than to blend
and almost identify the two, as a variety of circumstances has
occasioned or obliged us to do at this day.

It would at first sight seem
unnecessary to add to what has been said, any remark on mistakes or
apparent mistakes committed by S. Cyril in matters of fact; but as
this is often a ground of misconception, the subject shall be briefly
noticed. For instance, as to his statement concerning the discovery of
the True Cross \footnote{29. Lect.
iv. 10. x. 19. xiii. 4.}, he is
to be treated as any other historical witness under the same
circumstances, and the weight of his evidence, whatever it is, is to
be balanced against the improbability of the fact recorded, whether
antecedent, or arising from the silence concerning it of Eusebius and
Constantine. Again, we may well allow that he was not a natural
historian, without hurting his theological character. It is true that
he believed in the existence of the Phœnix \footnote{30. xviii.
8.}, and argued from the analogy afforded by it in favour of the
Resurrection. That is, he was philosophical on false grounds. And in
like manner persons have proved, as they thought, the Noachical deluge
from stones found on the top of hills, or have attributed it to the
action of a comet, or have believed or doubted the existence of
the sea-serpent or the dodo, and never have been reckoned worse or
better divines for either success of failure in such conjectures. It
as little follows that a theologian must be an ornithologist, as that
an ornithologist or a comparative anatomist must be a theologian; and
as no one in this day would reckon ignorance of divinity as a bar to
eminence and authority in scientific researches, so it betrays a
poverty of argument to reproach S. Cyril, or Eusebius, or S. Clement
before them, with not being proficients in a branch of knowledge which
has been a peculiar study of modern times. They did not profess to be
natural historians; let it be enough for this age to cultivate
physical science itself, without molesting the Fathers with its new
standards of intellectual superiority. Let it be enough for it to
despise the province of theology, without seeking to remodel it. The
Fathers did not profess the science on which it prides itself; nothing
but inspiration could secure them from shewing ignorance concerning
it; and no one pretends that S. Cyril or S. Clement were inspired.

It is only necessary to add
with respect to the present Translation, that for almost the whole of
it the Editors are indebted to Mr. CHURCH, Fellow of Oriel College. It has been
made from the Benedictine Text compared with the Oxford Edition of
Milles, the Benedictine Sections in the separate Lectures being marked
by numbers at the beginning of\emph{}the paragraphs, and the Oxford
sections on the margin. The few notes which are introduced are almost
confined to the elucidation of matter of fact, and have been kept
clear as far as possible from the expression of opinions; in drawing
them up, much use has been made of the valuable information contained
in the Oxford and Benedictine Editions. Such words of S. Cyril as have
a theological, controversial, or critical importance, are usually
placed in the margin opposite their place in the Translation.
The quotations from Scripture are given in the words of our received
version, wherever the Greek of Cyril admitted of it; when otherwise,
it has been signified in the margin.

J. H. N.

\emph{O\textsc{xford} The Feast of St. Matthew}, 1838.

\xchapter{The Treatises of S. Cæcilius Cyprian, Bishop of Carthage, and Martyr,}

\xsection{Translated, with notes and indices}

[Parker and Rivington, 1840]

\xsection{Preface}

THE
Treatises of St. Cyprian may suitably be preceded by the short Memoir
of his life written by his Deacon Pontius, and the Proconsular Acts of
his Martyrdom.

The Memoir is recommended to
our attention, not so much by any special excellence in itself, as by
the circumstance that it is written by one who was about the Bishop's
person, who attended him in exile, and who was a witness of his death
\footnote{1. S. Jerome (Script. De Vir.
Illust. 68.) praises this life as an ``egregium volumen.'' Ancient
Martyrologies record that Pontius eventually followed his master in
Martyrdom. The Bollandists, however, distinguish between him and the
Martyr Pontius, who was a Priest, and suffered in Piedmont.}. The reader need scarcely
be reminded, that the Deacon in St. Cyprian's age, as afterwards, was
the personal attendant and minister of the Bishop; thus St. Laurence
is celebrated as Deacon or Archdeacon to Sextus or Xystus, Bishop of
Rome and Martyr, the contemporary of St. Cyprian; and St. Athanasius
as Deacon to Alexander, Bishop of Alexandria, in the Council of Nicæa.

The Proconsular Acts are
considered to be the substance of the original, with the incidental
additions of subsequent times \footnote{2. The
substantial authenticity of these Acts seems to be generally allowed;
by the Benedictines, by Cave, Lit. Hist. art. Pontius, and by Gibbon,
who says that they and Pontius' life ``are consistent with each other
and with probability.'' The Bollandists consider that the Confession
and Martyrdom were ``extracted by the faithful from the public Acts,
and then a few words added in order to form them into a continuous
narration. And that in like manner some additions were made at the end
concerning the mode and circumstances of the Martyrdom, \&c.''}.

What has further to be said of
St. Cyprian is reserved for the second part of\emph{}the Volume,
which will contain his Letters. It shall only be added here, that he
was converted to the Christian faith\emph{}about A.D. 246, consecrated A.D.
248, and martyred A.D. 258.

\xsubsection{The life of St. Cyprian, by Pontius his Deacon}

CYPRIAN, that
religious Priest and glorious Witness of God, composed many works,
whereby may survive the memory of so worthy a name; the abundant
fecundity of his eloquence, and of God's grace in him, so widely
spread itself in copiousness and richness of speech, that perchance
even to the end of the world he will speak on; and yet, forasmuch as
his works and merits claim as a right that they should become an
example to us in writing, it has seemed good to draw up this brief
summary of it; not as if the life of so great a man were unknown\emph{}to
any of the heathen, but that even to our posterity may be handed on
his singular and high example unto an immortal memory. Certainly it
were hard, when even laymen and catechumens, who have obtained
martyrdom, have been honoured by our forefathers for their very
martyrdom's sake, with a record of many, nay of all details of their
passion, in order to our acquaintance with it who were yet unborn,
hard were it to pass over Cyprian's passion, so great a Priest and so
great a Martyr, who even over and above his martyrdom had lessons to
teach; and hard again to hide the deeds which he did in his life.
Those in truth were such, so great and wonderful, as to deter me by
the very contemplation of their greatness, and to urge me to a
confession of my incapacity to do justice to my subject, or to
represent his high deeds in correspondent terms, except that the
multitude of his achievements tells its own tale without heralding
from others. It has to be added, that you too are longing to hear
much, or, if possible, the whole concerning him, having a burning
desire at least to know his deeds, though his word of mouth be silent.
In which respect to say that I am deficient in the resources of eloquence, is to say little. Eloquence itself fails of the means of
fully satisfying your longing. Thus we are sorely pressed on either
side; by the weight of his excellences, by the importunity of your
entreaties.

A.D. 246.

2. From what shall I commence? where enter
upon his excellences, but from faith as a first principle, and from
his heavenly birth? considering that the deeds of a man of God should
be reckoned from no other point than that of his being born of God. He
might have employments before it, and a heart engaged and imbued with
liberal arts; still I pass over all this, as up to this date tending
merely to advantage of this life \footnote{3. S.
Gregory Nazianzen, in his oration in praise of S. Cyprian, (Orat. 18.)
states, that before his conversion he was addicted to magical arts,
which he made use of against a Christian female, named Justina, of
whom he was enamoured; that she however betook herself to Christ and
St. Mary, and the attempt ended in his burning his books, and
professing Christianity. Fell rejects the account altogether as a mere
fiction, (Monit. in Conf. S. Cypr.); Maranus, the Benedictine Editor,
(in vit.) and Tillemont refer it to a Cyprian, Bishop of Antioch in Phœnicia,
who has a place in both the Roman and Greek calendars. S. Cyprian was
a teacher of rhetoric, of great reputation; Jerom. de Vir. Illustr.
67. and before his conversion seems to have plunged into the usual
excesses of heathenism. vid. Treatise i. 2, 3. He seems not to have
been a native of Carthage. vid. Ep. 7. ed. Fell. St. Austin seems to
speak of him as a Senator. Serm. 311. c. 7.}. But after he had learned sacred knowledge and had emerged out
of the clouds of this world into the light of spiritual wisdom,
whatever I was witness to, whatever I have discovered of his
preferable works, I will relate; with the request that those
deficiencies of my narrative, which I feel will occur, should be
charged upon my ignorance rather than on his fame.

3. While he was yet in the rudiments of his
faith [i.e. before Baptism],
he felt that nothing was more fitting towards God than the observance
of continence; for the breast became what it should be, and the
understanding reached the full capacity of truth, when the lust of the
flesh was trampled on with the healthy and unimpaired vigour of
sanctity. Who has ever recorded such a marvel? the second birth had
not yet given eyes to the new man in the full radiance of divine
light, yet he was now conquering the old and previous darkness by the
mere outskirts of that light. Next, what is greater still, when he had
gained from Scripture certain lessons not according to the measure of
his noviciate but with the rapidity of faith, he at once appropriated to himself what he there read to be profitable in
meriting of the Lord. Diverting his property to the maintenance of the
indigent, and distributing whole estates in money, he secured two
benefits at once, both renouncing the pursuit of this world, than
which nothing is more pernicious, and observing mercy;—mercy, which
God has preferred even to His sacrifices, in which even he failed who
said that he had kept all the commandments of the law, and by which
with an anticipating haste of piety [vid. infra i. 3], he arrived at
perfection almost before he had learned how \footnote{4. S.
Cyprian himself attributes his change of heart and life to his
baptism; and while confessing with Pontius ``to sin no more has come of
faith,'' declares also, ``after that lifegiving water succoured me, what
was dark began to shine, what seemed impossible, now could be
achieved,'' i. 3.}. Who, let me ask, of the ancients, has done this? who of the
most esteemed elders in the faith, whose minds and ears have through
ever so many years been assailed by the words divine, ventured any
thing such as he, this man of an unformed faith and perchance
unrecognized profession, did achieve, surpassing the old time by
glorious and admirable works? No one reaps as soon as he has sowed.
None treads out the vintage from a young plantation. None yet ever
sought ripe fruit of bushes freshly planted. In him all things
incredible met together. In him the threshing anticipated, (if it can
be said, for the thing surpasses belief,) anticipated, I say, the
sowing; the vintage the tendril; the fruit, the firm root.

4. The Epistle of the Apostle
[1 Tim. 3, 6] says, that novices should be passed by; lest the
drowsiness of heathenism hanging on the scarcely rallied senses,
unlearned freshness might offend in aught against God. He was the
first, and, I suppose, the sole instance, that greater progress is
made by faith than by time. That Eunuch indeed in the Acts of the
Apostles is described as being baptized at once by Philip, because he
believed with his whole heart; but the parallel does not hold. For the
one was both a Jew, and in his way from the Lord's Temple was reading
the Prophet Isaiah, and had hope in Christ, though he thought Him not
yet come; the other, coming of the unlearned heathen, had as ripe a
faith at first, as few perhaps have at last. In a word, there was no
delay in his case as to the grace of God [i.e. Baptism], no
postponement. I have said too little: he forthwith received the
Presbyterate and Priesthood [A.D. 247]. Who
indeed would not commit all the ranks of honour to such a mind
believing? Many are the things he did when yet a layman, many when a
Presbyter, many after the example of just men of old, with a close
imitation, earning of the Lord, and surrendering himself to all the
duties of religion. And whenever he read of any one who had been
mentioned with praise by God, this was his ordinary advice, that we
should inquire on account of what deeds he had pleased God. If Job,
glorious by the testimony of God, is called a true worshipper of God,
one to whom no one might be compared on earth, he taught that ``one
ought to do whatever Job had done before; that, while we too do the
same, we may obtain the same testimony of God upon ourselves. Job,
despising the ruin of his estate, was so strong in practised virtue,
as not to feel even temporal losses of his benevolence. Penury broke
him not, nor grief, neither his wife's prayers, nor his bodily
sufferings shook his resolution. Virtue remained fixed in her own
home; and resignation established upon deep foundations, was moved by
no assault of the devil who tempted, from blessing his Lord with a
thankful faith even amid adversity. His house was open to any one who
came. No widow returned with her lap empty; nor blind, but was guided
by him as a companion; nor feeble in step, but was lifted by him as by
a carrier; nor helpless under the hand of the powerful, but had him
for a champion. ``These, things,'' he used to say, ``must they do
who would please God.'' \footnote{5. This
passage does not occur in any of S. Cyprian's extant Treatises; it
resembles them in style.}
And thus running through the specimens of all good men, while he ever
imitated the best, he set forth himself also for imitation.

5. He had an intimacy with one
among us, a just and memorable man, by name Cæcilius \footnote{6. S.
Cyprian, adopted as a Christian name, the name of the one to whom he
owed so much; vid. Jerom. l. c. Hence his full names are Thascius Cæcilius
Cyprianus.}, a Presbyter both by age and order, who had converted him from
his wanderings in this world to the acknowledgment of the true
divinity: him he loved with full honour and all observance, looking up
to him with dutiful veneration, not merely as the friend and brother
of his soul, but as though the parent of his new life. And so it
was that Cæcilius, comforted by such attentions, was led, and
reasonably, to such a fulness of affection, that, on departing from
this world, when his summons was near, he commended to him his wife
and children, and thus, from making him a member of his communion, in
the event made him the heir of his affection \footnote{7.
Clerics, however, ``by the Canons of the African Church, could not
become trustees to the property of their brethren, on the ground that
they were bound to serve nought but the altar and the sacrifice, and
to keep their time for supplications and prayers.'' Fell in Cypr. Epist.
1. vid. Conc. Carthag. A.D. 348. The same rule may be alluded to in
Treatise vi. 4. infr. ``Numerous Bishops, despising their sacred
calling, engaged themselves in secular vocations,'' ``divinâ
procuratione contemptâ, \emph{procuratores} rerum secularium fieri.''}. It were long to go through details; it were a toil to
enumerate his holy deeds.

A.D. 248.

6. For evidence of his good works, I suppose this is enough, that by
the judgment of God and the good will of the people, he was chosen for
the office of the Priesthood, and the rank of the Episcopate, while
yet a neophyte, and, as was considered, a novice \footnote{8. Vid. 1
Tim. iii. 6. S. Ambrose, Nectarius, Eusebius of Cæsarea in Cappadocia,
and others, were made Bishops under the same circumstances.}. Although still in the first days of his faith, and in the
rudimental season of his spiritual life, in such sort did his noble
disposition shine out, that, resplendent in the brightness at least of
hope, though not of office, he promised a full performance of the
duties of the priesthood, which was coming on him. Nor will I pass
over that special circumstance, how, while the whole people, God
influencing, poured itself out in love and honour of him, he on the
other hand humbly withdrew himself, yielding to older men, and deeming
himself unworthy of the title of such honour, whereby he became the
more worthy. For he is but made more worthy, who declines what he
deserves. With such emotion was the excited people at that time
agitated, longing with spiritual desire, as the event proves, not a
Bishop merely; but in him who had hid himself, and whom it was by a
divine presage so demanding, seeking, not a Priest only, but a Martyr
to come. A numerous brotherhood had beset the doors of his house;
solicitous love poured itself around all the approaches. What befel
the Apostle might then perhaps have been granted to him, as he wished
it, to be let down through a window; had he already shared with the
Apostle the honour of ordination. One might see all others in
anxious suspense waiting for his coming, and receiving him with excess
of joy when he came. I say it unwillingly, but I must say it. Some
resisted him \footnote{9. Five
priests opposed his consecration, one of them being Novatus; they
afterwards fomented the disorders of which the confessors were made
the instrument, (vid. infra \emph{Introd}, to Treatise v.) and joined
the patty of Felicissimus. This they did when S. Cyprian was in
concealment during the persecution. vid. Ep. 43. init. ed. Fell.}, even that
he might obtain his wish. Whom however, how forbearingly, how
patiently, how kindly he bore with! how indulgently he forgave,
reckoning them afterwards among his most intimate and familiar
friends, to the wonder of many! for who, but might count it miraculous
that so retentive a memory should become so oblivious?

7. How henceforth he bore
himself, who would suffice to relate! how great was his
loving-kindness, his strength of mind! his mercy, his severity! Such
sanctity and grace shone forth from his countenance as to confuse the
gazer. His look was grave and glad; neither a sternness which was sad,
nor overmuch good nature; but a just mixture of both; so that one
might doubt whether he claimed more our reverence or our love, except
that he claimed both. Nor did his dress belie his countenance,
subdued, as it was, to the middle course. He was not the man to be
inflated with the pride of the world's fashions; yet neither to grovel
in a studious penury; in that the latter style of dress is as
boastful, as that so ambitious frugality is ostentatious. How, when a
Bishop, he acted towards the poor, whom he already loved as a
catechumen, let the priests of mercifulness consider; whether taught
in the office of good works by the discipline of their very order, or
obliged to the duty of love by the general bond of the Gospel
Sacrament. As for Cyprian, what he was, such his Bishop's seat found
him ready made, and did not make him.

A.D. 250.

8. And so it was that for such merits he forthwith obtained also the
glory of proscription. Nor was it other than fitting that one, who
within the retreat of conscience so abounded in the full honours of
religion and faith, should also have a public name among the Gentiles.
Indeed he might even then, for the rapidity with which he developed
into all things, have hastened to the appointed crown of Martyrdom;
especially since the cries were frequent which called him ``to
the lion;'' \footnote{10. ``Christianos
ad Leonem.'' Tertullian Apol. 40. de Spect. 26.} had it not
been meet that he should pass through all degrees of glory before he
came to the highest, and had not the ruin of the Church which then
threatened needed the aid of so fertile a mind. For imagine him taken
hence at that time [vid. Treat. i.] by
the high reward of Martyrdom; who was there to shew the gains of grace
making progress by faith? who to curb the single women as it were with
the bridle of the Lord's lessons into a congruous rule of chastity [vid.
iv.], and a dress becoming their holiness [vid.
vi.]? who to teach penitence to the Lapsed [vid.
vii.]? truth to heretics, unity to
schismatics? to the sons of God peace and the law of Gospel prayer?
who to be the instrument of overthrowing blaspheming Gentiles, by
retorting on them their charges on us [vid.
ii. and viii.]? by whom were Christians,
grieved at loss of friends with excess of fondness or (what is worse)
defect of faith [vid. x.],
by whom to be comforted with the hope of things to come [vid. x.]?
from whom should we else learn mercy? from whom patience [vid. xi.]?
who was there to repress the evil feeling springing from the malignity
of poisonous envy [vid. xii.], with the sweetness of a salutary remedy
[vid. xiii.]? who to cheer the host of Martyrs with the exhortation of
a divine discourse,—who lastly to hasten with a stirring heavenly
trumpet those many confessors, signed with a second inscription on
their brow, and reserved as living examples of Martyrdom? Well surely
it was ordered then, well and indeed divinely, that a man so necessary
for so many and so good objects, was retarded from a Martyr's
consummation \footnote{11. S.
Jerome relates, that he had seen an old man, who professed to have
seen in his youth an amanuensis of S. Cyprian's, who was in the habit
of relating that the latter never passed a day without reading
Tertullian, continually saying to him, Da Magistrum; Hand me my
Master. vid. Jerom. de Vir. Illustr. 53. also Introd. to Treatise iv.
That S. Cyprian however did not follow Tertullian implicitly is plain
from his retiring from the persecution, not to mention other points of
difference.}.

9. You wish to be sure that
that retirement of his which now took place, was not from fear \footnote{12. ``On
the subject of flight in persecution, vid. infra note g, on vi. 8. vid.
also Ep. 34. fin. ed. Fell. Tertullian in his Montanistic Tract \emph{De
fugâ in} \emph{Persecutione} maintains that flight is unlawful.
The Roman Clergy (Ep. 8.) find fault with S. Cyprian's flight: he
defends himself, (Ep. 20.) saying he withdrew to hinder a riot. His
warrant for doing so was a divine direction. vid. Ep. 16. ``When a
persecution impended, the Bishops used to assemble the people, and
exhort them to constancy. Then they baptized infants and catechumens,
and divided the Eucharist among the faithful.'' Vales. In Euseb. Hist.
viii. 11. S. Dionysius was accused of having retired without first
attending to these necessary duties. \emph{ibid.}}; not to allege other evidence, he did suffer afterwards;
which suffering of course he would have shrunk from according to his
wont, had he shrunk from it before. But in truth, fear it was, but
right fear; fear of offending the Lord, fear which had rather be
dutiful to God's precepts, than be crowned together with the breach of
them. A mind surrendered in all things to God, and a faith enslaved to
the divine directions, considered that it would be sinning in very
suffering, unless it had obeyed the Lord who then ordered that
retreat. Something more must here be said on the advantage of the
postponement, though already I have touched on the subject. By what
seems shortly to have taken place, we may prove, as follows, that that
retirement did not issue from human pusillanimity, but, as is the
case, was really divine. The people of God had been ravaged with the
extraordinary and fierce assaults of a harrassing persecution; and,
whereas the crafty enemy could not deceive all by one and the same
artifice, therefore raging against them in manifold ways, wherever the
incautious soldier exposed his side, there he worsted each by various
overthrows. Some one was required who, when wounds had been received,
and darts cast by the changeful art of the torturing enemy, had
heavenly remedies at hand according to the nature of each, now to
pierce and now to sooth; and then was preserved a man of a mind beyond
all others divinely tempered, to steer the Church in a steady middle
course between the rebounding waves of colliding schisms. Let me ask
then, is not such design divine? could it have been without God's
governance? Let them look to it who think that such things happen by
chance. The Church answers to them with loud voice, declaring that she
does not allow, does not believe, that these her necessary champions
are reserved without the providence of God.

A.D. 252.

10. However, let me be allowed to run through the rest. A dreadful
pestilence broke out afterwards \footnote{13. For a
description of the pestilence, vid. infra ix. 9. vid. also the letters
of Dionysius of Alexandria (Euseb. Hist. vii. 22.) and S. Gregory
Nyssen's life of Gregory of Neo-Cæsarea, in fin. In the year 262 it
was especially destructive in Rome and in the cities of Greece,
carrying of in Rome as many as 5000 persons daily. Half the population
of Alexandria perished in it, according to Gibbon, who says that it ``raged
without interruption in every province, every city, and almost every
family of the Roman empire, from 250 to 265.'' Hist. x. fin. Its
duration is variously estimated.}, and the extraordinary ravages of a hateful sickness
entered house after house of the trembling populace in succession,
carrying off with sudden violence numberless people daily, each from
his own home. There was a general panic, flight, shrinking from the
infection, unnatural exposure of infected friends; as though to carry
the dying out of doors, were to rid one's self of death itself.
Meanwhile multitudes lay about the whole city, not bodies, but by this
time corpses; and called on the pity of passers-by from the view of a
fortune common to both parties. No one looked to aught beyond his
cruel gain. No one was alarmed from the recollection of parallel
instances. No one did to another what he wished done to himself. It
were a crime to pass over what in such circumstances was the conduct
of this Pontiff of Christ and God, who had surpassed the Pontiffs of
this world as much in benevolence as in truth of doctrine. First he
assembled the people in one place, urged on them the excellence of
mercifulness, taught them by instances from holy Scripture how much
the offices of benevolence avail to merit with God. Then he subjoined
that there was nothing wonderful in cherishing our own with the
fitting dutifulness of charity; that he became the perfect man, who
did somewhat more than publican or heathen, who, overcoming evil with
good and exercising what resembled a divine clemency, loved even his
enemies, who prayed, as the Lord admonishes and exhorts, for the
well-being of those who are persecuting him. He then makes His sun
rise, and bestows rain from time to time to foster the seed, shewing
forth all these benefits not only to His own, but to strangers also;
and he, who professes himself even God's son, why\emph{}follows he
not the example of his Father? ``We should answer to our birth,'' he
says; ``it is not fit that they should be degenerate who are known to
have been born again by God; rather the seed of a good Father should
be evidenced in the offspring, by our copying of His goodness.'' I pass
over many other things and those important, which my limits will not
allow me to detail; about which let it suffice to have noticed thus much. If the very Gentiles, had they heard them in the rostrum, would
probably have believed forthwith, what should a Christian people do,
whose very name begins in faith? Accordingly ministrations are divided
among them at once, according to the ranks and circumstances of such.
Many who from stress of poverty were unable to shew forth benefits of
cost, shewed forth what was more than costliness; by their personal
toil doing other services more precious than all riches. Who indeed
under such a teacher but must haste to be occupied in some part of
that warfare, by which he would be pleasing God the Father, and Christ
the Judge, and so good a Priest besides? Accordingly they did good in
the profusion of exuberant works to all, and not only to the household
of faith. They did somewhat more than is recorded of the incomparable
benevolence of Tobias. He must pardon the word, again pardon it,
pardon it often; or, to speak more truly, he must in equity grant,
that, although there was room for very much before Christ, yet after
Him there has been room for somewhat more, since to Christ's times the
fulness is ascribed. The slain of the king and the outcasts, whom
Tobias gathered together, were of his own kin only.

A.D. 257.

11. To these so good and so merciful deeds banishment succeeded. For
unbelief ever makes such return, recompensing the worse for the
better. Nor need I mention what God's Priest answered the proconsul
who questioned him, for there are Acts which relate it. Any how he is
forbidden the city, he who had done some good towards its health; he
who had toiled lest the eyes of the living should suffer the horrors
of the infernal abode; he, I say, who sleepless in the watchings of
benevolence had by a blameless kindness, (O the crime!) secured a
deserted state and destitute country from the sight of many exiles,
when all were flying from the loathsome look of the city. But this is
the world's concern in it, with whom exile is a punishment. To us our
country is less dear, who have a name in common, who abhor even our
own parents if they would persuade us contrary to the Lord. To them it
is a heavy punishment to live away from their city. To the Christian
the whole world is our home. Wherefore, though he be sent away into
ever so hidden and remote a place, having share in the things of
his God, he cannot count it banishment. Besides, while he serves God
entirely, even in his own city he is a stranger. For while he abstains
from desires of the flesh by continence of the Holy Ghost, putting off
the conversation of the old man, he is a foreigner even among his
citizens, or, I may say, among the very parents of his earthly life.
Moreover, though this might seem a punishment under other
circumstances, yet in such causes and sentences which we suffer for
trial of our virtue, it is not punishment, it is glory. But even
suppose banishment to be a punishment to us. If so, they are guilty of
the most extreme of crimes and the worst impiety, as their own
conscience testifies, who bring themselves to visit the innocent with
what they deem a punishment. I will not at present delineate a
delightful spot; I say nothing at first of the addition of all kinds
of beauties. Let us suppose the place offensive in its circumstances,
wretched to look upon, without wholesome water, or pleasant green, or
neighbouring shore; with vast rocks covered with forests, amid the
inhospitable depths of an altogether desert solitude, far off in the
world's trackless districts. Such a place might indeed bear the name
of exile, had Cyprian, priest of God, come thither; to whom if man's
ministrations failed, even the birds as to Elias, or the Angels as to
Daniel, would minister. Far, far indeed be it from any one to believe,
that even the least among us, provided he remained in the confession
of the Holy Name, should want any thing; so far was he God's Pontiff,
who had ever been urgent in matters of mercifulness, from wanting the
aid of all these things.

12. Next let us recount with
thanksgiving what I had put as the second supposition; namely, that
there was divinely provided for the soul of such a man, a sunny and
sufficient place, a place of sojourn, secret, as he could wish it, and
whatever has been before promised as his portion who seeks the kingdom
and righteousness of God \footnote{14.
Curubis, the place of S. Cyprian's exile, was ``a free and maritime
city of Zeugitania, in a pleasant situation, a fertile territory, and
at a distance of about forty miles from Carthage.'' Gibbon, Hist. ch.
16.}.
And, not to dwell upon the frequent visits of his brethren, nay, the
love of the very citizens, which afforded to him all things
whereof he seemed to be despoiled, I will not pass over the wonderful
visitation of God, by which He willed His Priest to be so sure in
exile of his passion which was to follow, that from his more abundant
assurance of the impending Martyrdom, Curubis possessed not an exile
only, but even a Martyr. For on that day when first we remained in the
place of banishment, (for me he chose out of his household in the
condescension of his love to be a voluntary exile, which, O had I been
also in his passion!) ``there appeared to me,'' said he, ``before I was
yet sunk in slumber, a young man greater than the human stature, by
whom being led as if to the prætorium, I seemed to myself to be
brought near to the tribunal of the proconsul then sitting. He, on
seeing me, forthwith began to write down upon a tablet a sentence,
which I knew not, for he had not asked me questions in the usual form;
however, that young man, who stood behind his back, with great anxiety
read whatever had been set down. And, since he could not utter it in
words, he intimated it by signs, which declared what was in the
writing of that tablet. For opening his hand and flattening it like a
blade, and imitating the blow of customary execution, he expressed
what he would have signified as if in clear words. I understood the
future sentence of my passion. I began at once to ask and seek, that
the delay even of one day might be given me, in order to my settling
my affairs in a regular way. After I had frequently repeated my
prayer, he began again to set down something on the tablet. I
perceived however, from the sereneness of his countenance, that the
judge's mind was influenced by the request, as if reasonable.
Moreover, that youth, who already had divulged somewhat by gesture, if
not by word, concerning my passion, made haste to signify by secret
signs from time to time, twisting his fingers one behind another, that
the delay was granted which I asked until the morrow. For me, although
the sentence was not read, while my heart exulted at the pleasant news
of delay granted, yet such was my alarm, from the chance of mistaking
the interpretation, that it was still all in flutter and agitation
from the remains of apprehension.''

13. What revelation could be
more manifest? what condescending mercy more blessed? All that
happened after in due course, were announced to him beforehand.
In nothing did the words of God come short; in nothing was the holy
promise mutilated. Do but review each particular as it was shewn to
him. He seeks a delay till the morrow when his sentence of suffering
was under deliberation; alleging his wish to settle his affairs on the
day which he had gained. His one day signified a year, which he was to
pass in this world after the vision. For, to speak more distinctly, he
was crowned, at the completion of the year, on that very day, on which
this had been announced to him at its commencement. For the day of the
Lord, though we do not find it used for year in divine Scripture, yet
in making promise of things to come, we consider that that space of
time ought to be given. Hence it matters not, if nothing short of a
year be announced while a day was spoken of, since that would
necessarily be more complete, which is greater. And whereas it was
explained by gesture not by speech, express speech was reserved for
the presence of the time itself. For it is usual then to set forth a
thing in words, when what is set forth is actually fulfilled. For no
one knew for certain wherefore this was shewn to him, till it turned
out that he was crowned on the same day on which he had seen it. And
yet in the interval his impending passion was known for certain by
all; but as to the particular day of his passion all those very
persons were silent, as if they were ignorant. And indeed I find some
such thing in the Scriptures. For the Priest Zacharias, when a son was
promised him by the Angel, because he believed not, became dumb; so
that by signs he asked for a tablet, seeing he had, not to utter, but
to write his son's name. Reasonably here too, when God's messenger
signified the Bishop's impending passion mainly by signs, he both
administered his faith and fortified his Priest. But again the reason
for seeking delay was his arranging his affairs and settling his will.
Now what affairs had he, what will to arrange, except Ecclesiastical
matters? For this reason there is a final delay granted, that
arrangements may be made as to whatever wants arrangement by a final
determination concerning the maintenance of the poor. And I consider
that for this sole end and for nothing else was he thus indulged by
those who had banished and were to kill him, that while here he
might relieve the poor who were here, with whatever remained to be
given of his final bounty, or, to speak more exactly, with the total
of his means. When then he had arranged matters so mercifully, and
thus ordered them in his last wishes, tomorrow's day drew near.

14. And now a messenger came
to him from the City from Xystus, that good and peace-making Priest,
and therefore most blessed Martyr [From Rome. Sextus. A.D. 258.]. The
executioner was expected every day, who was to strike through that
devoted neck of our most holy victim; and by this daily expectation of
dying, every day, as it came, became to him as though a day of
crowning. Meanwhile there came to him numbers of eminent and
illustrious persons, men of rank and family and secular distinction,
who, for the sake of their old friendship with him, urged him many
times to retire, backing their solicitations with the offer of
suitable places. But he, with mind hanging upon heaven, had put the
world out of sight, and did not assent to their persuasive
solicitations. Perhaps he would have done then also, what was urged on
him, and by many of the faithful too, if he had been bidden by divine
command \footnote{15. He
did at first retire and conceal himself at the advice of his friends.
This was on the Proconsul's coming to Utica; on the latter's returning
to Carthage, he came back to his gardens, and remained there, without
moving farther, till the officers arrested him. He had sold his
gardens on his conversion, but they had come back to him, perhaps (as
Gibbon supposes) by the kindness of his friends. vid. Pontius infr.
15. The opening of Treatise i. may stand for a description of them.}. Nor must we
leave unheralded the sublime glory of such a man, in that, when the
world was now raging and in reliance on its Rulers breathing out
hatred of the sacred Name, he, as occasion was given, fortified God's
servants with exhortations of the Lord, and animated them to tread
under foot sufferings of the present time, on the contemplation of the
glory which is to follow. In truth, there was in him so great a love
of sacred discourse, that while he prayed for passion, he desired that
it might be granted him while he was conversing concerning God.

15. And these were the daily
acts of a Priest destined for a sacrifice, pleasing to God; when
behold at the orders of the Proconsul, the Prætor's Official with his
soldiers suddenly surprised his gardens, those gardens which in the
beginnings of his faith he had sold, and, when God's kindness restored
them, would certainly have sold again for the benefit of the
poor, but that he feared to raise the jealousy of his persecutors. The
Official surprised him, or, I should more truly say, thought he had.
For what is there to surprise, as though by unforeseen attack, the
mind which is always ready? He went forward therefore, now certain
that that would be accomplished, which had long been held back; he
went forward with high and erect mind, with cheerfulness in his look,
and constancy in his heart. But being remanded till the morrow, he
turned from the Prætorium to the Official's house, when suddenly the
report spreads throughout Carthage, that ``Thascius was now brought
out,'' whom all knew, not only by the reputation in which he was
honourably held, but also from the recollection of his great
achievement [viz. in the plague]. All men throng together to a sight,
which for us was glorious from the self-sacrifice of his faith, but to
the Gentiles deplorable. However, during his lodgment for one night in
the house of the Official, his confinement was not rigorous, so that
we his intimates and friends were in his company as usual. Meanwhile
the whole people, conscious lest ought might be done in the night
without its own knowledge, kept watch at the door of the house. The
Divine goodness granted to him at that time, deserving as he was of
it, that God's people should even then keep vigil to usher in the day
of their Priest's Martyrdom. Some one, however, may perhaps ask, what
was the reason why he returned from the Prætorium to the Official;
and some think this, that on his part the Proconsul was then
unwilling. Far be it from me in things divinely overruled to complain
of indolence or caprice in the Proconsul. Far be it from me to allow
such an evil within the thoughts of a scrupulous mind, as that the
idle words of man should give sentence upon so blessed a Martyr. But
that next day, which a year before God's condescension had predicted,
was destined to be truly the morrow.

16. At length that other day
dawned, that appointed, promised, divine day \footnote{16. S.
Cyprian suffered on the same day as Cornelius of Rome, and six years
after him.}; which though the tyrant himself had desired to put off, he
would not at all have been able; a day pleasant in the secret
knowledge of the Martyr who was to be, all clouds being
dispersed throughout the world's circuit, and the sun shining
brightly. He left the Official's house, he an Official of Christ and
God, being hemmed in by the crowds of a mixed multitude on every side.
So infinite an army joined his train, it seemed as though he was
coming with troops in array to subdue death. As he went, he had to
pass the race-course. Well did it happen, and as if with a meaning,
that he should pass by the place of a corresponding contest, who was
running for the crown of righteousness, and had just finished his
labours. When he reached the Prætorium, the Proconsul not yet having
arrived, a private room was allowed him. There, while he sate
profusely perspiring after his long journey, (it so happened that his
seat was covered with linen \footnote{17. ``The
Bishop's seat used anciently to be covered with linen.'' Ed. Ben.},
as if to secure to him the honours of the episcopate even under the
very stroke of Martyrdom,) one of the officers \footnote{18.
tesserariis; those who communicated the \emph{tessera} through the
century.}, who was formerly a Christian, offered him clothes of his own;
thinking he might be willing to exchange his moist garments for his
own dry ones, and for himself ambitious of nothing further in return
for his gift, than to possess the now bloody sweat of the Martyr on
his road to God. But he made answer, ``That were seeking remedy for
discomforts, which perchance may not last out the day.'' Is it
surprising that he thought light of weariness in body, who in soul had
made light of death? But, to be brief, suddenly the Proconsul is
announced; and he is brought out, placed before him, asked his name;
he says who he is, and no more.

17. Upon this the judge reads
from the tablet the sentence, which before in the vision he had not
read; a divine sentence, not lightly to be spoken; a sentence worthy
of such a Bishop and such a Witness; a glorious sentence, in which he
is called a ``standard-bearer of the sect,'' and ``an enemy of the gods,''
and one who should be made ``an example to his followers,'' and whose
blood should now be shed ``in vindication of the law.'' [vid. infra
Procons. Acts.] Most satisfactory, most true is this sentence; for
every thing that was said, though said by a Gentile, is divine. Nor
surely is it wonderful, that High Priests are apt to prophesy of the
passion. He had been a standard-bearer, who was in the
practice of teaching concerning the bearing of Christ's cross; an
enemy of the gods, who bade destroy idols; he was an example to his
own, who unto the many who were about to follow in the same way, first
of his province \footnote{19. i.e.
in the province so called, the Eastern or Proconsular Africa.}
presented these first-fruits of Martyrdom. In his blood too ``the law
began to be ratified,'' but the law of Martyrs, who rivalling their
teacher in an initiation of a like glory, themselves too ratified the
law of his example in their own blood.

18. And when he passed out of
the doors of the Prætorium, a crowd of soldiers accompanied him, and
that nothing might be wanting in his passion, centurions and tribunes
were at his side. The place where he was to suffer is level,
surrounded with numerous trees so as to afford a sublime spectacle.
But, whereas its exceeding breadth hindered the view amid that
tumultuous crowd, persons who favoured him had climbed up the
branches, that he might gain this distinction also, (as in Zacchæus's
history,) of being seen from the trees. And now his eyes being bound
with his own hands, he tried to hasten the delay of the executioner,
whose business is the steel; and who with failing hand and trembling
fingers scarce could grasp it, until, when the hour was ripe for his
glorification, that centurion was granted strength to consummate the
death of a rare man, his hand being nerved with power from above. O
blessed people of the Church, who in eyes and other senses and in
uplifted voice, suffered together with such a Bishop, and thus, as
they had always heard him discourse, were crowned by God the Judge!
For although it could not happen, as the common wish was, that the
whole people at once should suffer in partnership of his glory, yet
whoever had the hearty will to suffer under the eyes of Christ and in
the ears of His Priest, did by the sufficient witness of his wish,
send up his name God-wards, as if by a representative. And thus, his
passion being consummated, it came about, that Cyprian, who had been
an example to all good men, was moreover the first in Africa to dye
his priestly diadems \footnote{20. i.e.
his crowns of sanctity and priesthood become a crown of martyrdom. The
Romanists would make such passages as this allude to the tonsure. The
African Bishops cut their hair in a circle. Vallars. in Hieron. Ep.
142. vid. also August. Ep. 33. §. 5. Bingham does not dissent; though
he is ``not confident that this was the reason of the name \emph{coronati.}''
Antiqu. vi. 4. §. 17.} in blood. For from the time that the Episcopal Order is catalogued in
Carthage, none is ever related, even of the holiest Priests, to have
attained unto passion \footnote{21. S.
Cyprian himself seems to say that African Bishops had already been
martyred. Ep. 66. ed. Fell. Accordingly, Tillemont suggested that
Pontius speaks only of Africa in a restricted sense, or the
Carthaginian territory, which was called especially ``the Province,''
vid, supra 17. Baronius, Lumper, and others interpret it of Carthage
only, referring to the words which follow in Pontius' text. Others
understand Pontius to speak only of the Valerian Persecution. Gibbon
eagerly seizes on Pontius' assertion in its broadest sense, and uses
it for his own purpose.},
though service devoted to God is always counted in dedicated men as if
a martyrdom. But Cyprian reached even unto the perfect crown the Lord
consummating; so that in that very city in which he had so lived, and
had been the first to do such noble deeds, he was the first also to
decorate the ensigns of the heavenly priesthood with glorious
bloodshed. What shall I here do? between joy at his passion, and grief
at bereavement, my mind is divided, and two sorts of feelings oppress
a breast too straitened for them. Shall I grieve that I was not his
companion? but his triumph is to be celebrated. Shall I celebrate his
triumph? but I am in grief that I am not his companion. To you,
however, the truth is to be avowed, and simply, as you know it, that
it was in my purpose to be so. In his glory I exult much and more than
much, and yet I grieve more that I remain behind.

\xsubsection{The Confession and Martyrdom of St. Cyprian, from the Proconsular Acts}

A.D. 257. Aug 30.

WHEN the Emperor
Valerian was Consul for the fourth, and Gallienus for the third time,
on the third of the Kalends of September, Paternus Proconsul at
Carthage in his council-chamber thus spoke to Cyprian the Bishop. 'The
most sacred Emperors Valerian and Gallienus have honoured me with
letters, wherein they enjoin that all those who use not the religion
of Rome, shall formally make profession of their return to the use of
Roman rites; I have made accordingly enquiry of your name; what answer
do you make to me?' Cyprian the Bishop spake, 'I am a Christian
and Bishop; I know no other Gods besides the One and true God, who
made heaven and earth, the sea, and all things therein; this God we
Christians serve, to Him we pray day and night, for ourselves, for all
mankind, for the health of the Emperors themselves.' \footnote{22. Vid.
in like manner Polycarp. ad Phil. 12. Just. M. Apol. l. i. 17. Athenag.
Leg. 37. Tertullian, Apol. 30. Origen, in Cels. viii. 73. Euseb. Hist.
vii. 11.} Paternus Proconsul said, 'Do you persist in this purpose?'
Cyprian Bishop answered, 'That good purpose, which hath once
acknowledged God, cannot be changed.' Paternus Proconsul said, 'Will
you then, obeying the mandate of the Emperors, depart into exile to
the city of Curubis?' Cyprian Bishop said, 'I go.' Paternus Proconsul
said, 'The letters, wherewith I have been honoured by the Emperors,
speak of Presbyters as well as of Bishops; I would know of you
therefore, who be they, who are Presbyters in this city?' Cyprian
Bishop answered, 'By your laws you have righteously and with great
benefit forbidden any to be informers \footnote{23. For
this law vid. Justinian Cod. x. 11.}; therefore they cannot be discovered and denounced by me; but
they will be found in their own cities.' Paternus Proconsul said, 'I
am accordingly inquisitor in this place.' Cyprian said, 'Our rules
forbid any man to offer himself for punishment, and your ordinances
discourage the same; they may not therefore offer themselves \footnote{24. Vid.
August. Contr. Gaudent. i. 40. (31.) where this passage is referred
to. vid. also Cypr. Ep. 81. ed. Fell. Those who studiously exposed
themselves to persecution were called \emph{Professors}. vid. Lumper
in Vit. Cypr.}, but they will be discovered by your inquisition.' Paternus
Proconsul said, 'They shall be discovered by me;' and added, 'they
further ordain, that no conventicles be held in any place, and that
the Christians shall not enter their cemeteries; if any transgress
this wholesome ordinance, it shall be capital.' Cyprian Bishop
answered, 'Do as you have been instructed.'

Thcn Paternus the Proconsul
bade them lead away the Bishop Cyprian into exile. During his long
abode in this place, Aspasius Paternus was succeeded by Galerius
Maximus, who bade the Bishop Cyprian be
recalled from exile, and brought before him [A.D.
258.]. Cyprian, the holy Martyr, chosen of God, returned from
Curubis, to which he had been exiled by order of Aspasius Paternus
then Proconsul, and by sacred command abode in his own gardens.
There he was in daily expectation that he should be visited as it had
been shewn him [vid. supra Life 12.]. While he dwelt there, suddenly
on the Ides of September [Sept. 13.], in the consulship of Tuscus and
Bassus, there came to him two chief officials \footnote{25.
Principes; they were the chief officers of the Prætorian court.}; one the chief gaoler \footnote{26.
Strator officii. al. stator vid. Ducange in verb.} in the Proconsular court of Galerius, the other \footnote{27.
Equistrator.} marshal of the guard in the same court; they placed him
between them in a chariot, and carried him to Sexti \footnote{28. Sexti,
as it is written by Tillemont and Lumper, was a place according to
some authorities six miles, according to others, four miles from
Carthage. Morcelli writes it Sextum.}, whither the Proconsul had retired for the recovery of his
health. By order of the Proconsul he was reserved for hearing on
another day; so the blessed Cyprian was privately lodged in the house
of the chief gaoler of the court of the most honourable \footnote{29.
Clarissimi. vid. Gibbon Hist. ch. 17, who says that in the reigns of
the Antonines this title was the ordinary and legal style of senators.
Afterwards it was given to the governors of provinces.} Galerius Maximus, Proconsul, in the street which is called
Saturn's, between the temples of Venus and of Salus. Thither flocked
the whole multitude of the brethren; which when holy Cyprian knew, he
bade that the young women should be protected, seeing they all
continued in the open street before the gate of the officer's house.
So on another day, the 18th of the Kalends of October [Sept.14.], a
great crowd was collected early at Sexti, as the Proconsul commanded.
And the same day Cyprian was brought before him as he sat for judgment
in the court called Sauciolum \footnote{30. i.e.
the criminal court. vid. Ducange, and Fell in loc.}. The Proconsul demanded, 'Are you Thascius Cyprianus?' Cyprian
Bishop answered, 'I am he.' Galerius Maximus Proconsul said, ``The most
sacred Emperors have commanded you to conform to the Roman rites.''
Cyprian Bishop said, ``I refuse to do so.'' Galerius: ``Take heed for
yourself.'' Cyprian; ``Execute the Emperor's orders; in a matter so
manifest I may not deliberate.'' Galerius, after briefly conferring
with his judicial council, with much reluctance pronounced the
following sentence. ``You have long lived an irreligious life [sacrilega
mente.], and have drawn together a number of men bound by an unlawful
association \footnote{31.
Nefariæ conspirationis. Christianity was not recognized as a \emph{religio
licita} till the next year, 259, by Gallienus. vid. Neander Hist.
(Rose) vol. i. Sect. i. 2. A.}, and
professed yourself an open enemy to the gods and the religion
of Rome; and the pious, most sacred, and august Emperors, Valerian and
Gallienus, and the most noble Cæsar Valerian, have endeavoured in
vain to bring you back to conformity with their religious
observances;—whereas then you have been apprehended as principal and
ringleader in these infamous crimes, you shall be made an example to
those whom you have wickedly associated with you: the authority of law
shall be ratified in your blood.'' He then read the sentence of the
court from a written tablet. ``It is the will of this court, that
Thascius Cyprianus be immediately beheaded.'' Cyprian Bishop maid, ``Thanks
be to God.'' \footnote{32. Vid.
S. Augustin. Serm. 309. §. 6. which in several points illustrates and
confirms this narrative.} After
sentence was pronounced, the whole assembled of the brethren cried
out, ``We will be beheaded with him.'' A great tumult arose among the
brethren, and a crowd followed to the place of execution. He was
brought forth into the field near Sexti, where having laid aside his
upper garment \footnote{33.
Lacerna or byrrus, a cloke, anciently, of a red colour. Ducange.
Baronius would interpret it of the episcopal dress of his day; but the
passage in the Acts is an addition. vid, Bingham Antiqu. vi. 4. §.
18.}, he
kneeled down, and addressed himself in prayer to the Lord. Then
stripping himself of his dalmatic, and giving it to the Deacons, he
stood in his linen tunic \footnote{34. The
tunicle or dalmatic ``was used in the earliest ages of the Christian
Church. Originally it had no sleeves …
It is said that wide sleeves were added … about the fourth century
in the West … The English Ritual directs it to be used by the
assistant ministers in the Holy Communion.'' Palmer's Origines.
Appendix §. 4.},
and awaited the executioner, to whom when he came Cyprian bade five
and twenty pieces of gold be given. The brethren meanwhile spread
linen cloths and napkins on the ground before him. Being unable to tie
the sleeve of his robe at the wrist, Julian Presbyter and Julian
Subdeacon performed this office for him. Then the blessed Cyprian
covered his eyes with his hands, and so suffered. His body was exposed
in a place hard by, to gratify the curiosity of the heathen. But in
the course of the night it was removed, and transported with prayers
and great pomp with wax tapers and funeral torches to the burying
ground of Macrobius Candidianus the Procurator, near the fish ponds in
the Mappalian Way. A few days after, Galerius Maximus the Proconsul
died.

Thus suffered the most blessed
Martyr Cyprian, on the eighteenth day of the Kalends of October [Sep.
14.], under Valerian and Gallienus Emperors;
in the kingdom of our Lord Jesus Christ, to whom be honour and glory
for ever and ever. Amen.

———————

Some such notice of St.
Cyprian's life and death, as the above, was necessary to introduce the
following Treatises; the force of which, as compositions, depends in
no small degree on some previous knowledge of the character and
history of the writer. They are the words of one who loved
Christianity well enough to give up for it at a mature age secular
engagements, settled habits and opinions, property, quiet, and at
length life itself. While exhorting to alms giving, he is already an
example of voluntary poverty; if he praises virginity, he has himself
embraced the single life; he insists on the nothingness of things
earthly, having first chosen contempt and reproach; he denounces the
heathen magistrate, with the knowledge that he is braving his power;
and he is severe with the Lapsed, because he himself is to be a
Martyr. Without going into the details of his theological and
ecclesiastical career, these facts are the great outlines of his
history, and may suitably and profitably be set against the subjects
treated in the following pages, and his mode of treating them. So much
is there of pretence in the world; so easy is it to see truths which
are hard to practise, so skilful is the intellect in simulating moral
greatness, so quick to feel and admire the truth, and so dexterous in
expressing and adorning it, that we naturally look out for some
assurance, which professions seldom supply, that we are reading what
is real and spontaneous, and not a mere\emph{}semblance of high
qualities.

As regards the Translation,
for almost the whole of which the Editors are indebted to the
Rev. CHARLES
THORNTON, of Christ Church, it need only be
stated, that neither the text of Baluzius nor of Fell has been
followed implicitly, but, where they differed, one or other has been
preferred according to the particular case. An attempt has been made,
in one portion of the Scripture references, to mark S. Cyprian's
variations from the present Vulgate version; but the differences
between the latter and his own, though often considerable, are often
so small, as to make it a matter of nice judgment when he should be
said to agree or disagree with it. It would seem on the whole that the
Vulgate and S. Cyprian's version differ from each other most in the
Prophets, next in the rest of the Old Testament, and least in the
Gospels and Epistles. The Psalms must be excepted from this
comparison, in which there is very little difference of translation at
all, perhaps from substitution of the Vulgate on the part of
transcribers. Next to the Psalms, there is least difference in the
books of the Apocrypha, and among these in Ecclesiasticus. This
information and other assistance while the Volume has been in the
press, have been kindly supplied by two friends of one of the Editors.

J. H. N.

\emph{Oxford},\emph{}\\
\emph{Feast of St. Mark}, 1839.

\xchapter{Commentary on the Epistle to the Galatians and Homilies on the Epistle to the Ephesians of S. John Chrysostom, Archbishop of Constantinople}

\xsection{Translated, with notes and indices}

[Parker and Rivington, 1840]

\xsection{Preface}

ST.
CHRYSOSTOM'S Commentary on the Epistle to the Galatians is continuous, according to
chapter and verse; instead of being arranged in Homilies with a Moral
or Practical application at their close, as in his exposition of other
Epistles. It was written at Antioch, as Montfaucon infers from a
reference which the Author makes, upon ch. i. v. 16. (p. 20,) to other
of his writings, which certainly were written about the same time in
that city. vid. Hom. de Mutat. Nom. tom. iii. p. 98. ed. Ben. The year
is uncertain, but seems not to have been earlier than A.D. 395.

The
Homilies on the Epistle to the Ephesians have been by some critics
assigned to his Episcopate at Constantinople, in consequence of
certain imperfections in their composition, which seemed to argue
absence of the comparative leisure which he enjoyed at Antioch. There
is a passage too in Homily xi. p. 231, 2, which certainly is very
apposite to the Author's circumstances in the court of Eudoxia. Yet
there are strong reasons for deciding that they too were delivered at
Antioch. S. Babylas and S. Julian, both Saints of Antioch, are
mentioned familiarly, the former in Homily ix. p. 205, the latter in
Homily xxi. pp. 342, 3. Monastic establishments in mountains in the
neighbourhood are spoken of in Homily vi. p. 165, and xiii. p. 248 \footnote{1. Vid.
also xxi. p. 338.}; and those near Antioch are famous in St. Chrysostom's history.
A schism too is alluded to in Homily xi. p. 230, as existing in the
community he was addressing, and that not about a question of
doctrine; circumstances which are accurately fulfilled in the
contemporary history of Antioch, and which are more or less noticed in
the Homilies on 1 Cor. which were certainly delivered at Antioch \footnote{2.
Vid. also Preface to Transl. of Homilies on 1 Cor. p. xiii.}. Moreover, he makes mention of the prevalence of superstitions,
Gentile and Jewish, among the people whom he was addressing, in Homily
vi. fin., p. 166. Hom. xii. fin. p. 240. which is a frequent ground of
complaint in his other writings against the Christians of Antioch; vid.
in Gal. p. 15; in 1 Cor. Hom. xii. §. 13, 14; in Col. Hom. viii.
fin.; contr. Jud. i. p. 386-8. Since Evagrius, the last Bishop of the
Latin succession in the schism, died in A.D.
392, those Homilies must have been composed before that date.

As to the
Translations, the Editors have been favoured with the former by a
friend who conceals his name; and with the latter, by the Rev. WILLIAM
JOHN COPELAND,
M.A. Fellow of Trinity College, Oxford.

J. H. N.
